\documentclass[11pt,letterpaper]{article}
\PassOptionsToPackage{hyphens}{url}
\usepackage[margin=.83in,top=.8in,bottom=.8in,headheight=15pt]{geometry}
\usepackage[T1]{fontenc}
\usepackage{lmodern,amsmath,amssymb,mathtools,bm,booktabs,array,graphicx,tikz}
\usetikzlibrary{arrows.meta,positioning,calc,decorations.pathreplacing,patterns}
\definecolor{ink}{HTML}{17344D}\definecolor{accent}{HTML}{087E8B}\definecolor{pale}{HTML}{EDF5F5}\definecolor{gold}{HTML}{AD6B19}
\providecommand{\dd}{\,\mathrm{d}}\providecommand{\R}{\mathbb{R}}
\providecommand{\E}{\mathbb{E}}\providecommand{\Var}{\operatorname{Var}}
\providecommand{\Cov}{\operatorname{Cov}}\providecommand{\diag}{\operatorname{diag}}
\providecommand{\norm}[1]{\left\lVert#1\right\rVert}
\providecommand{\abs}[1]{\left|#1\right|}
\usepackage{microtype,hyperref,fancyhdr,enumitem,amsthm,sectsty,needspace}
\allsectionsfont{\raggedright}
\hypersetup{colorlinks=true,linkcolor=ink,urlcolor=accent,pdfauthor={Zan Ahmad},pdftitle={Osmotic water transport and cell volume / Review}}
\setlength{\parindent}{0pt}\setlength{\parskip}{6pt}\linespread{1.07}
\setlist{itemsep=3pt,topsep=4pt,leftmargin=*}
\emergencystretch=2em\clubpenalty=10000\widowpenalty=10000
\pagestyle{fancy}\fancyhf{}\fancyhead[L]{\small Mathematical Modeling}\fancyhead[R]{\small Lecture 8}
\fancyfoot[L]{\footnotesize Zan Ahmad}\fancyfoot[R]{\footnotesize\thepage}\renewcommand{\headrulewidth}{.3pt}
\newtheorem{proposition}{Proposition}\newtheorem{theorem}{Theorem}\newtheorem{lemma}{Lemma}
\theoremstyle{definition}\newtheorem{definition}{Definition}\newtheorem{example}{Example}
\numberwithin{equation}{section}
\begin{document}
\section*{Osmotic water transport and cell volume --- review sheet}
Use this after attempting the worked examples and practice. Explain each claim in words before calculating.

\section*{Key ideas}
\begin{itemize}
\item Amount n, volume V and concentration n/V are distinct.
\item Under equal hydrostatic pressure, ideal net water flow points toward higher impermeant-particle concentration.
\item With positive L, equilibrium volume is n/c\_out; conductance sets response speed.
\item The external concentration is fixed because the outside is modeled as a large reservoir.
\item Conservation, numerical accuracy and biological validation require different checks.
\end{itemize}
\section*{Common errors and repairs}
\textbf{Assuming concentration cannot change unless solute crosses.} Water movement changes V and therefore n/V even with fixed n.

\textbf{Using the initial flow unchanged for every later step.} Recompute inside concentration and flow from the current volume.

\textbf{Reading a stable numerical volume as proof of cell survival.} The model has no mechanics, rupture or viability process.

\section*{Retrieve without your notes}
\begin{itemize}
\item For n =\allowbreak{} 300 fmol and c\_out =\allowbreak{} 200 mM, what is V\_eq when L > 0?
\item If volume halves with fixed amount, what happens to concentration?
\item Why must the special case L =\allowbreak{} 0 be treated separately?
\end{itemize}
\section*{Readiness check}
\begin{itemize}
\item Distinguish trapped solute amount, water volume and concentration.
\item Predict inward or outward water flow under the stated osmotic assumptions.
\item Compute a short volume update with consistent units.
\item Derive an equilibrium volume and separate it from the speed of approach.
\item Check conserved solute, positive volume and numerical refinement without inferring cell survival.
\end{itemize}
\end{document}
